
\section{Implementation}

\subsection*{Problem Background}

In this project, we implement an optimized CUDA version of the forward pass of Andrej Karpathy's \textit{microgpt} \cite{karpathyMicrogpt2026}. \textit{Microgpt} is a minimal, dependency-free Python implementation of a GPT-style model. It is based on the GPT-2 architecture \cite{radfordLanguageModelsAre2019}, with a few changes made for simplicity. Both GPT-2 and Karpathy's \textit{microgpt} implement:
\begin{itemize}[noitemsep, topsep=0pt]
  \item Decoder-only Transformer architecture
  \item Token + positional embeddings
  \item Stacked Transformer blocks
  \item Causal (masked) multi-head self-attention
  \item Residual connections around attention and MLP sublayers
  \item Feedforward MLP with expansion and projection
  \item Autoregressive next-token prediction
\end{itemize}
Since \textit{microgpt} was intended to showcase the GPT-2 architecture in a single unoptimized Python file, the proposed model has only 4,192 parameters compared to GPT-2's 1.6 billion parameters. In addition to the reduced model size, \textit{microgpt} introduces several simplifications, including the use of RMSNorm instead of LayerNorm, removal of bias terms, and ReLU in place of GELU \cite{karpathyMicrogpt2026}. 

\noindent In order to keep the scope of our implementation to 10 hours, we focus only on speeding up the forward pass of \textit{microgpt}. Since we believed that comparing a CUDA C++ version to a Python version would introduce speedups due to both our parallelization choices and language differences (C++ vs Python), we first reimplemented \textit{microgpt} in C++, preserving the following:
\begin{itemize}[noitemsep, topsep=0pt]  \item Identical layer structure
  \item Same ordering of operations (norm $\rightarrow$ attention $\rightarrow$ residual $\rightarrow$ norm $\rightarrow$ MLP $\rightarrow$ residual)
  \item Same KV caching mechanism
  \item Same attention computation (scaled dot-product + softmax)
  \item Same MLP structure (\texttt{fc1} $\rightarrow$ activation $\rightarrow$ \texttt{fc2})
\end{itemize}

Our C++ version is nearly identical, with differences limited to implementation details such as replacing dynamic Python data structures with pre-allocated caches, removing autograd in favor of raw floating-point computation, and expressing all operations as explicit loop-based computations over contiguous memory.

\subsection*{C++ Serial Baseline}
Our implementation follows the training setup used in \textit{microgpt}, which operates on a character-level dataset consisting of a set of names $\mathcal{D} = \{s^{(1)}, s^{(2)}, \dots, s^{(N)}\}$. Each name $s^{(k)} \in \mathcal{D}$ is treated as a sequence of tokens, $s^{(k)} = [t_1, t_2, \dots, t_T]$, where $T$ denotes the sequence length. The model is trained autoregressively to predict the next token given all previous tokens. Our C++ implementation preserves this setup exactly, performing the same forward-pass computation and next-token prediction task.  
 
Given a tokenized sequence (i.e., a single name) $[t_1, t_2, \dots, t_T]$, the model computes a sequence of logits $[y_1, y_2, \dots, y_{T-1}]$, where each $y_i$ represents the model's prediction of token $t_{i+1}$ conditioned on the prefix $[t_1, \dots, t_i]$. In the baseline implementation, this computation is performed sequentially, iterating over the dataset one name at a time and, within each name, one token at a time, with each position attending over all previous tokens in the sequence.  A high level pseudocode of our serial implentation is shown under Algorithm~\ref{alg:forward-pass}.

We denote the embedding dimension as $d$, the number of Transformer layers as $L$, and the number of attention heads as $H$. For each token position $i$, the model produces a hidden representation $x_i \in \mathbb{R}^d$, which is transformed through attention and MLP operations across all layers. 

The computational cost of the forward pass over $N$ sequences is $O\!\left(N \cdot L \cdot (C_{\text{attn}} + C_{\text{mlp}})\right)$, where $C_{\text{attn}} = O(T^2 \cdot d)$ and $C_{\text{mlp}} = O(T \cdot d^2)$. Since the names have approximately constant sequence lengths, we can simplify this to $O(N \cdot L \cdot d^2)$. 


\begin{algorithm}[H]
\caption{High-Level Forward Pass}
\label{alg:forward-pass}
\begin{algorithmic}
\State $\text{loss} \gets 0$
\For{each sequence $s^{(k)} \in \mathcal{D}$}
    \For{$i = 1$ to $T-1$}
        \State $x_i \gets \text{TokenEmbedding}(t_i) + \text{PositionalEmbedding}(i)$
        \State $x_i \gets \text{RMSNorm}(x_i)$

        \For{$\ell = 1$ to $L$}
            \State $x_i \gets x_i + \text{Attention}(\text{RMSNorm}(x_i), \{k_j, v_j\}_{j \le i})$
            \State $x_i \gets x_i + \text{MLP}(\text{RMSNorm}(x_i))$
        \EndFor

        \State $y_i \gets \text{Linear}(x_i)$
        \State $\text{loss} \gets \text{loss} - \log(\text{Softmax}(y_i)[t_{i+1}])$
    \EndFor
\EndFor
\State $\text{loss} \gets \text{loss} / \text{total\_tokens}$
\end{algorithmic}
\end{algorithm}
